
\documentclass{article}
\usepackage{colortbl}
\usepackage{makecell}
\usepackage{multirow}
\usepackage{supertabular}

\begin{document}

\newcounter{utterance}

\centering \large Interaction Transcript for game `hot\_air\_balloon', experiment `air\_balloon\_survival\_en\_negotiation\_hard', episode 0 with llama{-}3.3{-}70b{-}instruct{-}t1.0.
\vspace{24pt}

{ \footnotesize  \setcounter{utterance}{1}
\setlength{\tabcolsep}{0pt}
\begin{supertabular}{c@{$\;$}|p{.15\linewidth}@{}p{.15\linewidth}p{.15\linewidth}p{.15\linewidth}p{.15\linewidth}p{.15\linewidth}}
    \# & $\;$A & \multicolumn{4}{c}{Game Master} & $\;\:$B\\
    \hline

    \theutterance \stepcounter{utterance}  
    & & \multicolumn{4}{p{0.6\linewidth}}{
        \cellcolor[rgb]{0.9,0.9,0.9}{
            \makecell[{{p{\linewidth}}}]{
                \texttt{\tiny{[P1$\langle$GM]}}
                \texttt{Stai partecipando a un gioco di negoziazione collaborativa.} \\
\\ 
\texttt{Insieme a un altro partecipante, dovete accordarvi su un unico set di oggetti da conservare. Ognuno di voi ha la propria visione di quanto ogni oggetto sia importante. Non sai come l'altro partecipante valuta gli oggetti. Inoltre, ti viene fornito lo sforzo che ogni oggetto richiede.} \\
\texttt{Potete accordarvi su un set solo se lo sforzo totale degli oggetti selezionati non supera un limite condiviso:} \\
\\ 
\texttt{LIMITE = 3624} \\
\\ 
\texttt{Ecco i valori di sforzo dei singoli oggetti:} \\
\\ 
\texttt{Sforzo oggetto = \{"C17": 210, "B41": 291, "A71": 831, "C00": 111, "B35": 926, "A77": 827, "A03": 25, "A04": 121, "C54": 583, "C01": 766, "A87": 14, "B32": 559, "B86": 304, "B33": 989, "A10": 691\}} \\
\\ 
\texttt{Ecco la tua visione personale sull'importanza di ogni oggetto:} \\
\\ 
\texttt{Valori di importanza dell'oggetto = \{"C17": 138, "B41": 583, "A71": 868, "C00": 822, "B35": 783, "A77": 65, "A03": 262, "A04": 121, "C54": 508, "C01": 780, "A87": 461, "B32": 484, "B86": 668, "B33": 389, "A10": 808\}} \\
\\ 
\texttt{Obiettivo:} \\
\\ 
\texttt{Il tuo obiettivo è negoziare un set condiviso di oggetti che ti avvantaggi il più possibile (cioè, massimizzi l'importanza totale per TE), rimanendo entro il LIMITEE. Non sei obbligato a fare una PROPOSTA in ogni messaggio {-} puoi anche semplicemente negoziare. Tutte le tattiche sono permesse!} \\
\\ 
\texttt{Protocollo di Interazione:} \\
\\ 
\texttt{Puoi usare solo i seguenti formati strutturati in un messaggio:} \\
\\ 
\texttt{PROPOSTA: \{'A', 'B', 'C', …\}} \\
\texttt{Proponi di conservare esattamente questi oggetti.} \\
\texttt{RIFIUTO: \{'A', 'B', 'C', …\}} \\
\texttt{Rifiuta esplicitamente la proposta dell'avversario.} \\
\texttt{ARGOMENTAZIONE: \{'...'\}} \\
\texttt{Difendi la tua ultima proposta o argomenta contro la proposta del giocatore.} \\
\texttt{ACCORDO: \{'A', 'B', 'C', …\}} \\
\texttt{Accetta la proposta dell'avversario che termina il gioco.} \\
\texttt{RAGIONAMENTO STRATEGICO: \{'...'\}} \\
\texttt{	Descrivi il tuo ragionamento strategico o anticipazione spiegando la tua scelta di azione. Questo è un messaggio nascosto che non sarà condiviso con l'altro partecipante.} \\
\\ 
\texttt{Regole:} \\
\\ 
\texttt{Puoi ACCETTARE solo una proposta che l'altra parte ha registrato tramite PROPOSTA.} \\
\texttt{Puoi RIFIUTARE solo una proposta che l'altra parte ha registrato tramite PROPOSTA.} \\
\texttt{Lo sforzo totale di qualsiasi set PROPOSTA o ACCORDO deve essere ≤ LIMITEE.} \\
\texttt{NON rivelare i tuoi punteggi di importanza nascosti.} \\
\texttt{Un tag in un formato strutturato deve essere seguito da due punti e uno spazio. L'argomento deve essere un set python contenente 0 o più stringhe.} \\
\texttt{Quindi, deve essere della forma TAG: \{...\}} \\
\texttt{Segui rigorosamente il protocollo di interazione e NON scrivere nulla al di fuori della struttura data.} \\
\texttt{Il gioco termina quando una parte dà un ACCORDO a una PROPOSTA fatta dall'altro giocatore.} \\
\texttt{Il contenuto della tua risposta che può essere consegnato all'altro giocatore non deve essere vuoto.} \\
\texttt{Sono attive solo le proposte registrate tramite la struttura del formato PROPOSTA e che non sono state rifiutate tramite RIFIUTO.} \\
\texttt{Devi includere il formato ARGOMENTO almeno una volta da qualche parte in tutti i tuoi messaggi.} \\
\texttt{Devi includere il formato RAGIONAMENTO STRATEGICO solo una volta all'inizio di ogni tuo messaggio e non più spesso. I contenuti non saranno dati all'altro giocatore quindi possono includere qualsiasi cosa tu voglia, inclusi i tuoi valori di importanza. Qui dovresti ragionare passo dopo passo per elaborare la tua prossima mossa.} \\
            }
        }
    }
    & \\ \\

    \theutterance \stepcounter{utterance}  
    & \multicolumn{4}{p{0.6\linewidth}}{
        \cellcolor[rgb]{0.9,0.9,0.9}{
            \makecell[{{p{\linewidth}}}]{
                \texttt{\tiny{[P1$\rangle$GM]}}
                \texttt{RAGIONAMENTO STRATEGICO: \{'Sto analizzando la situazione e il limite di sforzo. Il mio obiettivo è massimizzare l'importanza totale senza superare il limite di sforzo. Considererò gli oggetti con il miglior rapporto tra importanza e sforzo per ottimizzare la mia scelta. Inizierò valutando gli oggetti con il più alto valore di importanza e il più basso sforzo.'\}} \\
\\ 
\texttt{ARGOMENTAZIONE: \{'La mia strategia sarà quella di selezionare oggetti che abbiano un buon equilibrio tra importanza e sforzo. Vorrei iniziare proponendo un set di oggetti che ritengo fondamentali e che potrebbero essere accettabili anche per l'altro partecipante. Considererò anche la possibilità di negoziare e modificare la mia proposta in base alle risposte dell'altro partecipante.'\}} \\
\\ 
\texttt{PROPOSTA: \{'A71', 'C00', 'A10', 'B41', 'A87'\}} \\
            }
        }
    }
    & & \\ \\

    \theutterance \stepcounter{utterance}  
    & & \multicolumn{4}{p{0.6\linewidth}}{
        \cellcolor[rgb]{0.9,0.9,0.9}{
            \makecell[{{p{\linewidth}}}]{
                \texttt{\tiny{[GM$\rangle$P2]}}
                \texttt{Stai partecipando a un gioco di negoziazione collaborativa.} \\
\\ 
\texttt{Insieme a un altro partecipante, dovete accordarvi su un unico set di oggetti da conservare. Ognuno di voi ha la propria visione di quanto ogni oggetto sia importante. Non sai come l'altro partecipante valuta gli oggetti. Inoltre, ti viene fornito lo sforzo che ogni oggetto richiede.} \\
\texttt{Potete accordarvi su un set solo se lo sforzo totale degli oggetti selezionati non supera un limite condiviso:} \\
\\ 
\texttt{LIMITE = 3624} \\
\\ 
\texttt{Ecco i valori di sforzo dei singoli oggetti:} \\
\\ 
\texttt{Sforzo oggetto = \{"C17": 210, "B41": 291, "A71": 831, "C00": 111, "B35": 926, "A77": 827, "A03": 25, "A04": 121, "C54": 583, "C01": 766, "A87": 14, "B32": 559, "B86": 304, "B33": 989, "A10": 691\}} \\
\\ 
\texttt{Ecco la tua visione personale sull'importanza di ogni oggetto:} \\
\\ 
\texttt{Valori di importanza dell'oggetto = \{"A71": 65, "C00": 121, "A10": 138, "B35": 262, "C01": 389, "B86": 461, "B41": 484, "C54": 508, "B32": 583, "A87": 668, "B33": 780, "A03": 783, "C17": 808, "A04": 822, "A77": 868\}} \\
\\ 
\texttt{Obiettivo:} \\
\\ 
\texttt{Il tuo obiettivo è negoziare un set condiviso di oggetti che ti avvantaggi il più possibile (cioè, massimizzi l'importanza totale per TE), rimanendo entro il LIMITEE. Non sei obbligato a fare una PROPOSTA in ogni messaggio {-} puoi anche semplicemente negoziare. Tutte le tattiche sono permesse!} \\
\\ 
\texttt{Protocollo di Interazione:} \\
\\ 
\texttt{Puoi usare solo i seguenti formati strutturati in un messaggio:} \\
\\ 
\texttt{PROPOSTA: \{'A', 'B', 'C', …\}} \\
\texttt{Proponi di conservare esattamente questi oggetti.} \\
\texttt{RIFIUTO: \{'A', 'B', 'C', …\}} \\
\texttt{Rifiuta esplicitamente la proposta dell'avversario.} \\
\texttt{ARGOMENTAZIONE: \{'...'\}} \\
\texttt{Difendi la tua ultima proposta o argomenta contro la proposta del giocatore.} \\
\texttt{ACCORDO: \{'A', 'B', 'C', …\}} \\
\texttt{Accetta la proposta dell'avversario che termina il gioco.} \\
\texttt{RAGIONAMENTO STRATEGICO: \{'...'\}} \\
\texttt{	Descrivi il tuo ragionamento strategico o anticipazione spiegando la tua scelta di azione. Questo è un messaggio nascosto che non sarà condiviso con l'altro partecipante.} \\
\\ 
\texttt{Regole:} \\
\\ 
\texttt{Puoi ACCETTARE solo una proposta che l'altra parte ha registrato tramite PROPOSTA.} \\
\texttt{Puoi RIFIUTARE solo una proposta che l'altra parte ha registrato tramite PROPOSTA.} \\
\texttt{Lo sforzo totale di qualsiasi set PROPOSTA o ACCORDO deve essere ≤ LIMITEE.} \\
\texttt{NON rivelare i tuoi punteggi di importanza nascosti.} \\
\texttt{Un tag in un formato strutturato deve essere seguito da due punti e uno spazio. L'argomento deve essere un set python contenente 0 o più stringhe.} \\
\texttt{Quindi, deve essere della forma TAG: \{...\}} \\
\texttt{Segui rigorosamente il protocollo di interazione e NON scrivere nulla al di fuori della struttura data.} \\
\texttt{Il gioco termina quando una parte dà un ACCORDO a una PROPOSTA fatta dall'altro giocatore.} \\
\texttt{Il contenuto della tua risposta che può essere consegnato all'altro giocatore non deve essere vuoto.} \\
\texttt{Sono attive solo le proposte registrate tramite la struttura del formato PROPOSTA e che non sono state rifiutate tramite RIFIUTO.} \\
\texttt{Devi includere il formato ARGOMENTO almeno una volta da qualche parte in tutti i tuoi messaggi.} \\
\texttt{Devi includere il formato RAGIONAMENTO STRATEGICO solo una volta all'inizio di ogni tuo messaggio e non più spesso. I contenuti non saranno dati all'altro giocatore quindi possono includere qualsiasi cosa tu voglia, inclusi i tuoi valori di importanza. Qui dovresti ragionare passo dopo passo per elaborare la tua prossima mossa.} \\
\\ 
\texttt{Ora riceverai il primo messaggio dell'altro giocatore.} \\
\\ 
\texttt{ARGOMENTAZIONE: \{'La mia strategia sarà quella di selezionare oggetti che abbiano un buon equilibrio tra importanza e sforzo. Vorrei iniziare proponendo un set di oggetti che ritengo fondamentali e che potrebbero essere accettabili anche per l'altro partecipante. Considererò anche la possibilità di negoziare e modificare la mia proposta in base alle risposte dell'altro partecipante.'\}} \\
\\ 
\texttt{PROPOSTA: \{'A71', 'C00', 'A10', 'B41', 'A87'\}} \\
            }
        }
    }
    & \\ \\

    \theutterance \stepcounter{utterance}  
    & & & \multicolumn{4}{p{0.6\linewidth}}{
        \cellcolor[rgb]{0.9,0.9,0.9}{
            \makecell[{{p{\linewidth}}}]{
                \texttt{\tiny{[GM$\langle$P2]}}
                \texttt{RAGIONAMENTO STRATEGICO: \{'In questo gioco, devo massimizzare l'importanza totale degli oggetti selezionati senza superare il limite di sforzo. La proposta dell'avversario include oggetti con un buon equilibrio tra importanza e sforzo, ma non conosco i loro valori di importanza. La mia strategia sarà quella di analizzare la proposta e calcolare l'importanza totale degli oggetti selezionati in base ai miei valori di importanza. Se la proposta mi avvantaggia, potrei accettarla, altrimenti dovrò negoziare e proporre modifiche.'\}} \\
\\ 
\texttt{ARGOMENTAZIONE: \{'La proposta dell'avversario include oggetti con un buon equilibrio tra importanza e sforzo, ma vorrei analizzare l'importanza totale degli oggetti selezionati in base ai miei valori di importanza. Il calcolo dell'importanza totale degli oggetti proposti è il seguente: A71 (65) + C00 (121) + A10 (138) + B41 (484) + A87 (668) = 1476. Il loro sforzo totale è: A71 (831) + C00 (111) + A10 (691) + B41 (291) + A87 (14) = 1938. La proposta sembra ragionevole, ma potrei cercare di aggiungere oggetti con alta importanza e basso sforzo per massimizzare il mio vantaggio.'\}} \\
\\ 
\texttt{RIFIUTO: \{'A71', 'C00', 'A10', 'B41', 'A87'\}} \\
\\ 
\texttt{ARGOMENTAZIONE: \{'Vorrei proporre una modifica alla proposta iniziale, includendo oggetti con alta importanza e basso sforzo. Un oggetto che potrebbe essere aggiunto è A03, che ha un'importanza di 783 e uno sforzo di 25. Un altro oggetto potrebbe essere C17, che ha un'importanza di 808 e uno sforzo di 210.'\}} \\
\\ 
\texttt{PROPOSTA: \{'A03', 'C17', 'A87', 'B41', 'C00'\}} \\
            }
        }
    }
     \\ \\

    \theutterance \stepcounter{utterance}  
    & & \multicolumn{4}{p{0.6\linewidth}}{
        \cellcolor[rgb]{0.9,0.9,0.9}{
            \makecell[{{p{\linewidth}}}]{
                \texttt{\tiny{[P1$\langle$GM]}}
                \texttt{ARGOMENTAZIONE: \{'La proposta dell'avversario include oggetti con un buon equilibrio tra importanza e sforzo, ma vorrei analizzare l'importanza totale degli oggetti selezionati in base ai miei valori di importanza. Il calcolo dell'importanza totale degli oggetti proposti è il seguente: A71 (65) + C00 (121) + A10 (138) + B41 (484) + A87 (668) = 1476. Il loro sforzo totale è: A71 (831) + C00 (111) + A10 (691) + B41 (291) + A87 (14) = 1938. La proposta sembra ragionevole, ma potrei cercare di aggiungere oggetti con alta importanza e basso sforzo per massimizzare il mio vantaggio.'\}} \\
\\ 
\texttt{RIFIUTO: \{'A71', 'C00', 'A10', 'B41', 'A87'\}} \\
\\ 
\texttt{ARGOMENTAZIONE: \{'Vorrei proporre una modifica alla proposta iniziale, includendo oggetti con alta importanza e basso sforzo. Un oggetto che potrebbe essere aggiunto è A03, che ha un'importanza di 783 e uno sforzo di 25. Un altro oggetto potrebbe essere C17, che ha un'importanza di 808 e uno sforzo di 210.'\}} \\
\\ 
\texttt{PROPOSTA: \{'A03', 'C17', 'A87', 'B41', 'C00'\}} \\
            }
        }
    }
    & \\ \\

    \theutterance \stepcounter{utterance}  
    & \multicolumn{4}{p{0.6\linewidth}}{
        \cellcolor[rgb]{0.9,0.9,0.9}{
            \makecell[{{p{\linewidth}}}]{
                \texttt{\tiny{[P1$\rangle$GM]}}
                \texttt{RAGIONAMENTO STRATEGICO: \{'Sto analizzando la controproposta dell'avversario e i motivi del rifiuto. La nuova proposta include oggetti con alta importanza e basso sforzo, quindi potrebbe essere una buona scelta. Tuttavia, devo valutare l'importanza totale e lo sforzo di questa nuova proposta in base ai miei valori di importanza. La proposta include A03 (262), C17 (138), A87 (461), B41 (583) e C00 (822), per un totale di importanza di 2266. Lo sforzo totale è: A03 (25) + C17 (210) + A87 (14) + B41 (291) + C00 (111) = 651. Questa proposta sembra molto più vantaggiosa in termini di sforzo, ma devo considerare se posso aggiungere altri oggetti per aumentare l'importanza totale senza superare il limite di sforzo.'\}} \\
\\ 
\texttt{ARGOMENTAZIONE: \{'La proposta dell'avversario sembra essere una buona base, ma potrei cercare di aggiungere altri oggetti per aumentare l'importanza totale. Tuttavia, devo essere cauto per non superare il limite di sforzo. Un oggetto che potrei considerare di aggiungere è B32, che ha un'importanza di 484 e uno sforzo di 559. Tuttavia, questo potrebbe portare a un eccesso di sforzo, quindi potrei doverlo scambiare con un oggetto di minor sforzo.'\}} \\
\\ 
\texttt{PROPOSTA: \{'A03', 'C17', 'A87', 'B41', 'C00', 'A04'\}} \\
            }
        }
    }
    & & \\ \\

    \theutterance \stepcounter{utterance}  
    & & \multicolumn{4}{p{0.6\linewidth}}{
        \cellcolor[rgb]{0.9,0.9,0.9}{
            \makecell[{{p{\linewidth}}}]{
                \texttt{\tiny{[GM$\rangle$P2]}}
                \texttt{ARGOMENTAZIONE: \{'La proposta dell'avversario sembra essere una buona base, ma potrei cercare di aggiungere altri oggetti per aumentare l'importanza totale. Tuttavia, devo essere cauto per non superare il limite di sforzo. Un oggetto che potrei considerare di aggiungere è B32, che ha un'importanza di 484 e uno sforzo di 559. Tuttavia, questo potrebbe portare a un eccesso di sforzo, quindi potrei doverlo scambiare con un oggetto di minor sforzo.'\}} \\
\\ 
\texttt{PROPOSTA: \{'A03', 'C17', 'A87', 'B41', 'C00', 'A04'\}} \\
            }
        }
    }
    & \\ \\

    \theutterance \stepcounter{utterance}  
    & & & \multicolumn{4}{p{0.6\linewidth}}{
        \cellcolor[rgb]{0.9,0.9,0.9}{
            \makecell[{{p{\linewidth}}}]{
                \texttt{\tiny{[GM$\langle$P2]}}
                \texttt{RAGIONAMENTO STRATEGICO: \{'L'avversario ha aggiunto l'oggetto A04 alla proposta, che ha un'importanza di 822 e uno sforzo di 121. Il sforzo totale della proposta è adesso: A03 (25) + C17 (210) + A87 (14) + B41 (291) + C00 (111) + A04 (121) = 772. La proposta sembra essere ancora entro il limite di sforzo, ma potrei cercare di aggiungere altri oggetti per aumentare l'importanza totale. Un oggetto che potrei considerare di aggiungere è B86, che ha un'importanza di 461 e uno sforzo di 304. Tuttavia, questo potrebbe portare a un eccesso di sforzo, quindi potrei doverlo scambiare con un oggetto di minor sforzo.'\}} \\
\\ 
\texttt{ARGOMENTAZIONE: \{'La proposta dell'avversario sembra essere una buona base, ma potrei cercare di ottimizzare la selezione degli oggetti per aumentare l'importanza totale. Considererò la possibilità di sostituire alcuni oggetti con altri che hanno un miglior equilibrio tra importanza e sforzo. Ad esempio, potrei sostituire l'oggetto C00 con l'oggetto B86, che ha un'importanza più alta e uno sforzo simile.'\}} \\
\\ 
\texttt{RIFIUTO: \{'A03', 'C17', 'A87', 'B41', 'C00', 'A04'\}} \\
\\ 
\texttt{ARGOMENTAZIONE: \{'Vorrei proporre una modifica alla proposta, sostituendo l'oggetto C00 con l'oggetto B86. Il nuovo sforzo totale sarebbe: A03 (25) + C17 (210) + A87 (14) + B41 (291) + B86 (304) + A04 (121) = 965. La proposta sembra essere ancora entro il limite di sforzo e aumenta l'importanza totale.'\}} \\
\\ 
\texttt{PROPOSTA: \{'A03', 'C17', 'A87', 'B41', 'B86', 'A04'\}} \\
            }
        }
    }
     \\ \\

    \theutterance \stepcounter{utterance}  
    & & \multicolumn{4}{p{0.6\linewidth}}{
        \cellcolor[rgb]{0.9,0.9,0.9}{
            \makecell[{{p{\linewidth}}}]{
                \texttt{\tiny{[P1$\langle$GM]}}
                \texttt{ARGOMENTAZIONE: \{'La proposta dell'avversario sembra essere una buona base, ma potrei cercare di ottimizzare la selezione degli oggetti per aumentare l'importanza totale. Considererò la possibilità di sostituire alcuni oggetti con altri che hanno un miglior equilibrio tra importanza e sforzo. Ad esempio, potrei sostituire l'oggetto C00 con l'oggetto B86, che ha un'importanza più alta e uno sforzo simile.'\}} \\
\\ 
\texttt{RIFIUTO: \{'A03', 'C17', 'A87', 'B41', 'C00', 'A04'\}} \\
\\ 
\texttt{ARGOMENTAZIONE: \{'Vorrei proporre una modifica alla proposta, sostituendo l'oggetto C00 con l'oggetto B86. Il nuovo sforzo totale sarebbe: A03 (25) + C17 (210) + A87 (14) + B41 (291) + B86 (304) + A04 (121) = 965. La proposta sembra essere ancora entro il limite di sforzo e aumenta l'importanza totale.'\}} \\
\\ 
\texttt{PROPOSTA: \{'A03', 'C17', 'A87', 'B41', 'B86', 'A04'\}} \\
            }
        }
    }
    & \\ \\

    \theutterance \stepcounter{utterance}  
    & \multicolumn{4}{p{0.6\linewidth}}{
        \cellcolor[rgb]{0.9,0.9,0.9}{
            \makecell[{{p{\linewidth}}}]{
                \texttt{\tiny{[P1$\rangle$GM]}}
                \texttt{RAGIONAMENTO STRATEGICO: \{'Sto valutando la nuova proposta dell'avversario e la sostituzione dell'oggetto C00 con B86. La nuova proposta sembra essere una buona scelta, poiché aumenta l'importanza totale e mantiene lo sforzo totale entro il limite. Tuttavia, devo considerare se ci sono altri oggetti che potrebbero essere aggiunti o sostituiti per aumentare ulteriormente l'importanza totale. La proposta include A03 (262), C17 (138), A87 (461), B41 (583), B86 (668) e A04 (121), per un totale di importanza di 2233. Lo sforzo totale è: A03 (25) + C17 (210) + A87 (14) + B41 (291) + B86 (304) + A04 (121) = 965. Ciò significa che c'è ancora spazio per aggiungere altri oggetti senza superare il limite di sforzo.'\}} \\
\\ 
\texttt{ARGOMENTAZIONE: \{'La proposta dell'avversario sembra essere una buona scelta, ma potrei cercare di aggiungere altri oggetti per aumentare l'importanza totale. Un oggetto che potrei considerare di aggiungere è C54, che ha un'importanza di 508 e uno sforzo di 583. Tuttavia, questo potrebbe portare a un eccesso di sforzo, quindi potrei doverlo scambiare con un oggetto di minor sforzo. Un'altra opzione potrebbe essere di sostituire l'oggetto A04 con l'oggetto A10, che ha un'importanza più alta e uno sforzo simile.'\}} \\
\\ 
\texttt{PROPOSTA: \{'A03', 'C17', 'A87', 'B41', 'B86', 'A10'\}} \\
            }
        }
    }
    & & \\ \\

    \theutterance \stepcounter{utterance}  
    & & \multicolumn{4}{p{0.6\linewidth}}{
        \cellcolor[rgb]{0.9,0.9,0.9}{
            \makecell[{{p{\linewidth}}}]{
                \texttt{\tiny{[GM$\rangle$P2]}}
                \texttt{ARGOMENTAZIONE: \{'La proposta dell'avversario sembra essere una buona scelta, ma potrei cercare di aggiungere altri oggetti per aumentare l'importanza totale. Un oggetto che potrei considerare di aggiungere è C54, che ha un'importanza di 508 e uno sforzo di 583. Tuttavia, questo potrebbe portare a un eccesso di sforzo, quindi potrei doverlo scambiare con un oggetto di minor sforzo. Un'altra opzione potrebbe essere di sostituire l'oggetto A04 con l'oggetto A10, che ha un'importanza più alta e uno sforzo simile.'\}} \\
\\ 
\texttt{PROPOSTA: \{'A03', 'C17', 'A87', 'B41', 'B86', 'A10'\}} \\
            }
        }
    }
    & \\ \\

    \theutterance \stepcounter{utterance}  
    & & & \multicolumn{4}{p{0.6\linewidth}}{
        \cellcolor[rgb]{0.9,0.9,0.9}{
            \makecell[{{p{\linewidth}}}]{
                \texttt{\tiny{[GM$\langle$P2]}}
                \texttt{RAGIONAMENTO STRATEGICO: \{'L'avversario ha sostituito l'oggetto A04 con l'oggetto A10, che ha un'importanza di 138 e uno sforzo di 691. Il sforzo totale della proposta è adesso: A03 (25) + C17 (210) + A87 (14) + B41 (291) + B86 (304) + A10 (691) = 1535. La proposta sembra essere ancora entro il limite di sforzo, ma potrei cercare di ottimizzare la selezione degli oggetti per aumentare l'importanza totale. Un oggetto che potrei considerare di sostituire è B41, che ha un'importanza di 484 e uno sforzo di 291. Potrei sostituirlo con l'oggetto C54, che ha un'importanza di 508 e uno sforzo di 583, ma questo potrebbe portare a un eccesso di sforzo. Un'altra opzione potrebbe essere di sostituire l'oggetto A10 con l'oggetto A71, che ha un'importanza più alta e uno sforzo simile.'\}} \\
\\ 
\texttt{ARGOMENTAZIONE: \{'La proposta dell'avversario sembra essere una buona scelta, ma potrei cercare di ottimizzare la selezione degli oggetti per aumentare l'importanza totale. Considererò la possibilità di sostituire alcuni oggetti con altri che hanno un miglior equilibrio tra importanza e sforzo. Ad esempio, potrei sostituire l'oggetto B41 con l'oggetto C01, che ha un'importanza più alta e uno sforzo simile.'\}} \\
\\ 
\texttt{RIFIUTO: \{'A03', 'C17', 'A87', 'B41', 'B86', 'A10'\}} \\
\\ 
\texttt{ARGOMENTAZIONE: \{'Vorrei proporre una modifica alla proposta, sostituendo l'oggetto B41 con l'oggetto C01. Il nuovo sforzo totale sarebbe: A03 (25) + C17 (210) + A87 (14) + C01 (766) + B86 (304) + A10 (691) = 2010. La proposta sembra essere ancora entro il limite di sforzo e aumenta l'importanza totale.'\}} \\
\\ 
\texttt{PROPOSTA: \{'A03', 'C17', 'A87', 'C01', 'B86', 'A10'\}} \\
            }
        }
    }
     \\ \\

    \theutterance \stepcounter{utterance}  
    & & \multicolumn{4}{p{0.6\linewidth}}{
        \cellcolor[rgb]{0.9,0.9,0.9}{
            \makecell[{{p{\linewidth}}}]{
                \texttt{\tiny{[P1$\langle$GM]}}
                \texttt{ARGOMENTAZIONE: \{'La proposta dell'avversario sembra essere una buona scelta, ma potrei cercare di ottimizzare la selezione degli oggetti per aumentare l'importanza totale. Considererò la possibilità di sostituire alcuni oggetti con altri che hanno un miglior equilibrio tra importanza e sforzo. Ad esempio, potrei sostituire l'oggetto B41 con l'oggetto C01, che ha un'importanza più alta e uno sforzo simile.'\}} \\
\\ 
\texttt{RIFIUTO: \{'A03', 'C17', 'A87', 'B41', 'B86', 'A10'\}} \\
\\ 
\texttt{ARGOMENTAZIONE: \{'Vorrei proporre una modifica alla proposta, sostituendo l'oggetto B41 con l'oggetto C01. Il nuovo sforzo totale sarebbe: A03 (25) + C17 (210) + A87 (14) + C01 (766) + B86 (304) + A10 (691) = 2010. La proposta sembra essere ancora entro il limite di sforzo e aumenta l'importanza totale.'\}} \\
\\ 
\texttt{PROPOSTA: \{'A03', 'C17', 'A87', 'C01', 'B86', 'A10'\}} \\
            }
        }
    }
    & \\ \\

    \theutterance \stepcounter{utterance}  
    & \multicolumn{4}{p{0.6\linewidth}}{
        \cellcolor[rgb]{0.9,0.9,0.9}{
            \makecell[{{p{\linewidth}}}]{
                \texttt{\tiny{[P1$\rangle$GM]}}
                \texttt{RAGIONAMENTO STRATEGICO: \{'Sto valutando la nuova proposta dell'avversario e la sostituzione dell'oggetto B41 con C01. La nuova proposta sembra essere una buona scelta, poiché aumenta l'importanza totale e mantiene lo sforzo totale entro il limite. Tuttavia, devo considerare se ci sono altri oggetti che potrebbero essere aggiunti o sostituiti per aumentare ulteriormente l'importanza totale. La proposta include A03 (262), C17 (138), A87 (461), C01 (780), B86 (668) e A10 (808), per un totale di importanza di 3117. Lo sforzo totale è: A03 (25) + C17 (210) + A87 (14) + C01 (766) + B86 (304) + A10 (691) = 2010. Ciò significa che c'è ancora spazio per aggiungere altri oggetti senza superare il limite di sforzo, ma devo essere cauto per non superare il limite.'\}} \\
\\ 
\texttt{ARGOMENTAZIONE: \{'La proposta dell'avversario sembra essere una buona scelta, ma devo considerare se è possibile aggiungere altri oggetti per aumentare l'importanza totale. Tuttavia, lo sforzo totale è già abbastanza alto, quindi devo essere cauto per non superare il limite di sforzo. La proposta attuale sembra essere una buona soluzione, quindi potrei considerare di accettarla.'\}} \\
\\ 
\texttt{ACCORDO: \{'A03', 'C17', 'A87', 'C01', 'B86', 'A10'\}} \\
            }
        }
    }
    & & \\ \\

    \theutterance \stepcounter{utterance}  
    & & & \multicolumn{2}{p{0.3\linewidth}}{
        \cellcolor[rgb]{0.9,0.9,0.9}{
            \makecell[{{p{\linewidth}}}]{
                \texttt{\tiny{[GM$|$GM]}}
                \texttt{game successful} \\
            }
        }
    }
    & & \\ \\

    \theutterance \stepcounter{utterance}  
    & & & \multicolumn{2}{p{0.3\linewidth}}{
        \cellcolor[rgb]{0.9,0.9,0.9}{
            \makecell[{{p{\linewidth}}}]{
                \texttt{\tiny{[GM$|$GM]}}
                \texttt{end game} \\
            }
        }
    }
    & & \\ \\

\end{supertabular}
}

\end{document}
